%% HalalChain — versi Bahasa Indonesia
%% Build: pdflatex main-indo; bibtex main-indo; pdflatex main-indo; pdflatex main-indo

\documentclass[journal]{IEEEtran}

\IEEEoverridecommandlockouts

\usepackage[indonesian]{babel}
\usepackage{cite}
\usepackage{amsmath,amssymb,amsfonts}
\usepackage{algorithmic}
\usepackage{graphicx}
\usepackage{textcomp}
\usepackage{xcolor}
\usepackage{booktabs}
\usepackage{array}
\usepackage{tabularx}
\usepackage{makecell}
\usepackage{placeins}
\usepackage{fancyvrb}
\usepackage{fvextra}
\usepackage{xurl}
\usepackage{tikz}
\usetikzlibrary{arrows.meta,positioning,shapes.geometric,calc}
\usepackage[hidelinks]{hyperref}

\graphicspath{{figures-indo/}}

\DefineVerbatimEnvironment{Code}{Verbatim}{%
  fontsize=\scriptsize,
  breaklines=true,
  breakanywhere=true,
  xleftmargin=0pt,
  xrightmargin=0pt
}

\begin{document}

\title{HalalChain: Protokol Ketertelusuran Batch Halal pada Base L2 dengan Integrasi Persyaratan Desain Etis Islam dan Jejak Audit Permanen}

\author{%
  \IEEEauthorblockN{Muhammad Fatih Maulana\IEEEauthorrefmark{1},
  Rio Ferdana Sudrajat\IEEEauthorrefmark{1},
  Muhammad Rahardian Baihaqi\IEEEauthorrefmark{1},\\
  Muhammad Fahriza Pratama\IEEEauthorrefmark{1},
  Yazid Zinky Arisona\IEEEauthorrefmark{1},
  dan Muhammad Fauzan Lubada\IEEEauthorrefmark{1}}
  \IEEEauthorblockA{\IEEEauthorrefmark{1}%
  Universitas Islam Negeri Sunan Gunung Djati Bandung, Bandung, Indonesia}
}

\markboth{HalalChain: Ketertelusuran Batch pada Base L2}{HalalChain: Ketertelusuran Batch pada Base L2}

\maketitle

\begin{abstract}
Ekonomi halal global melampaui USD~2 triliun dalam belanja konsumen, namun kepercayaan terhadap sertifikasi masih rapuh: sertifikat fisik dapat dipalsukan, basis data perusahaan dapat diubah oleh administrator yang berwenang, dan konsumen jarang memperoleh bukti tingkat batch pada saat pembelian.
Regulasi pelabelan halal wajib di Indonesia (BPJPH/MUI) memperkuat kebutuhan ketertelusuran yang dapat diverifikasi dan berbiaya rendah bagi usaha mikro, kecil, dan menengah (UMKM) yang tidak mampu memakai platform blockchain kelas enterprise.
Makalah ini memperkenalkan \emph{HalalChain}, protokol ketertelusuran batch terdesentralisasi yang di-deploy pada Base Network---rollup Ethereum Layer~2---yang mengaitkan \emph{status} sertifikasi, pengenal konten (CID), dan metadata audit pada ledger publik, sementara dokumen pendukung (PDF, laporan laboratorium, daftar bahan) disimpan pada InterPlanetary File System (IPFS).
Kontribusi kami meliputi: (i)~desain kontrak pintar minimal tiga peran dengan kontrol akses berbasis peran (RBAC) OpenZeppelin dan riwayat revisi permanen setelah penolakan auditor; (ii)~pemetaan eksplisit dari persyaratan \emph{desain etis} selaras Maqasid ke properti protokol yang dapat ditegakkan; (iii)~implementasi sumber terbuka yang terdiri atas kontrak Solidity, pengujian Hardhat, dan dasbor web produsen/auditor/konsumen dwibahasa dengan verifikasi QR tanpa dompet kripto bagi konsumen.
Bagian~\ref{sec:eval} melaporkan analisis keamanan kualitatif dan pengujian skenario fungsional; pada testnet Ethereum Sepolia, \texttt{registerBatch} menggunakan 155.506 gas (0,00078367 ETH) dan \texttt{verifyBatch} menggunakan 79.214 gas (0,00034832 ETH).
HalalChain menunjukkan bahwa integritas batch halal dapat diaudit secara publik tanpa mewajibkan konsumen mengoperasikan dompet kripto.
Prototipe langsung tersedia di \url{https://web-six-ivory-36.vercel.app/}.
\end{abstract}

\begin{IEEEkeywords}
Blockchain, ketertelusuran halal, rantai pasok, IPFS, Base Network, kontrak pintar, penskalaan Layer~2, etika Islam, RBAC
\end{IEEEkeywords}

\section{Pendahuluan}
\label{sec:intro}

Sertifikasi halal di Indonesia diatur oleh Badan Penyelenggara Jaminan Produk Halal (BPJPH) berdasarkan Undang-Undang No.~33/2014, dengan audit kategori produk yang sering melibatkan Majelis Ulama Indonesia (MUI).
Meskipun kemajuan regulasi, konsumen dan pasar ekspor masih melaporkan kesenjangan kepercayaan: sertifikat palsu, rantai pasok yang tidak transparan, dan keterbatasan kemampuan menghubungkan produk ritel dengan batch produksi tertentu yang diaudit~\cite{sidarto2021halal,rashid2018halal}.
Kekhawatiran serupa muncul dalam perdagangan halal lintas batas, di mana negara pengimpor tidak dapat secara mandiri memastikan bahwa label mencerminkan batch yang sama yang lulus sertifikasi domestik~\cite{marianingsih2024halal,rejeb2018halal}.

\subsection{Pernyataan Masalah}
Tiga kelemahan struktural mendorong penelitian ini.
\textbf{Pertama}, sertifikat berbasis kertas atau PDF bersifat \emph{dapat disalin}---dokumen palsu dapat tampak autentik tanpa akses ke registri resmi.
\textbf{Kedua}, basis data ketertelusuran terpusat---termasuk platform khusus halal dan sistem pangan generik---mengonsentrasikan wewenang tulis pada administrator yang pada prinsipnya dapat mengubah catatan historis~\cite{kshetri2018blockchain,underwood2016blockchain,viriyasitavat2019blockchain}.
\textbf{Ketiga}, percontohan ketertelusuran blockchain publik (misalnya IBM Food Trust, VeChain) menargetkan rantai pasok enterprise dan umumnya memerlukan infrastruktur dompet, sehingga tidak terjangkau bagi produsen UMKM yang mendominasi sektor pangan halal Indonesia~\cite{kamath2018blockchain,kamilaris2019blockchain}.

Teknologi blockchain menawarkan catatan kriptografis yang diamankan dan bersifat tambah-saja yang mendukung persyaratan desain \emph{Amanah} (kepercayaan)---kesaksian permanen dan tahan manipulasi tentang siapa menyetujui apa dan kapan~\cite{zheng2020overview,christidis2016blockchain}.
Namun, men-deploy dokumen lengkap pada Ethereum mainnet tidak ekonomis untuk pendaftaran batch volume tinggi~\cite{benet2014ipfs,rejeb2020blockchain}.
HalalChain mengatasi ketegangan ini dengan mengaitkan status minimal on-chain pada Base L2~\cite{base2024docs,optimism2024spec} sambil menyimpan bukti lengkap off-chain pada IPFS, serta mengodekan penolakan dan revisi auditor sebagai status protokol tingkat pertama, bukan koreksi di luar ledger.

\subsection{Pertanyaan Penelitian}
\begin{enumerate}
  \item \textbf{PP1 (Biaya dan kelayakan):} Apakah arsitektur kontrak pintar minimal dapat menyediakan verifikasi batch halal permanen pada tingkat biaya L2 yang terjangkau bagi produsen UMKM?
  \item \textbf{PP2 (Pertukaran arsitektur):} Bagaimana pemisahan on-chain/off-chain (IPFS) memengaruhi biaya gas, jaminan ketidakberubahan, dan ketersediaan dokumen?
  \item \textbf{PP3 (Keselarasan etis):} Persyaratan desain etis Islam mana yang dapat dipetakan ke properti sistem yang dapat diverifikasi dan diuji tanpa menyamakan perilaku perangkat lunak dengan putusan agama?
\end{enumerate}

\subsection{Kontribusi}
Makalah ini memberikan kontribusi sebagai berikut:
\begin{enumerate}
  \item \textbf{Desain protokol HalalChain}---model tiga peran (Produsen, Auditor, Konsumen) dengan RBAC OpenZeppelin~\cite{openzeppelin2024access}, status siklus batch eksplisit, dan riwayat revisi on-chain yang menghubungkan pengajuan ulang ke batch induk yang ditolak.
  \item \textbf{Analisis arsitektur}---kajian pertukaran antara jangkar on-chain (CID, status, cap waktu, alamat auditor) dan dokumen off-chain pada IPFS~\cite{benet2014ipfs,caro2018blockchain}, termasuk jalur baca konsumen tanpa dompet.
  \item \textbf{Pemetaan desain etis}---ketertelusuran dari persyaratan selaras Maqasid (Bagian~\ref{sec:prelim}) ke penegakan kontrak pintar, dibingkai sebagai persyaratan rekayasa, bukan fatwa.
  \item \textbf{Implementasi terbuka dan kerangka evaluasi}---basis kode yang dapat direproduksi dengan tujuh pengujian unit, tiga skenario fungsional, dan validasi keamanan kualitatif (Bagian~\ref{sec:eval-a}); tolok ukur empiris testnet Ethereum Sepolia (Bagian~\ref{sec:eval-b}) dan validasi skenario fungsional (Bagian~\ref{sec:eval-a}).
\end{enumerate}

\subsection{Organisasi Makalah}
Bagian~\ref{sec:related} meninjau literatur blockchain rantai pasok dan ketertelusuran halal.
Bagian~\ref{sec:prelim} menetapkan konteks sertifikasi dan persyaratan desain etis.
Bagian~\ref{sec:design}--\ref{sec:implementation} menjelaskan arsitektur, spesifikasi kontrak, dan implementasi web.
Bagian~\ref{sec:eval} menyajikan evaluasi; Bagian~\ref{sec:discussion} membahas keterbatasan; Bagian~\ref{sec:conclusion} menyimpulkan.

\section{Kajian Terkait}
\label{sec:related}

Tabel~\ref{tab:related} membandingkan HalalChain dengan sistem ketertelusuran representatif dalam hal spesifisitas halal, desentralisasi, kebutuhan dompet konsumen, dan dukungan riwayat revisi pasca-penolakan.
Platform khusus halal sebelumnya seperti HALAL TRACE mengandalkan basis data terpusat dan tidak mengekspos jejak audit publik yang permanen~\cite{rashid2018halal}.
Sistem ketertelusuran pangan blockchain umum (IBM Food Trust, VeChain) menunjukkan adopsi industri tetapi tidak disesuaikan dengan alur kerja sertifikasi halal, semantik penolakan auditor, atau pembingkaian etis Islam~\cite{kamilaris2019blockchain,casino2021systematic,kamath2018blockchain}.
Portal pemerintah seperti BPJPH SIHALAL menyediakan registrasi resmi tetapi tetap dikelola secara terpusat tanpa verifikasi publik tingkat batch bagi pihak ketiga mana pun.

\begin{table}[t]
  \caption{Perbandingan sistem ketertelusuran yang relevan dengan HalalChain.}
  \label{tab:related}
  \centering
  \footnotesize
  \setlength{\tabcolsep}{3pt}
  \begin{tabularx}{\columnwidth}{@{}l *{4}{>{\centering\arraybackslash}X}@{}}
    \toprule
    \textbf{Sistem} & \textbf{Halal} & \textbf{Desent.} & \textbf{Tanpa Dompet} & \textbf{Revisi} \\
    \midrule
    IBM Food Trust & Tidak & Parsial & Tidak & Tidak \\
    VeChain & Tidak & Ya & Tidak & Tidak \\
    HALAL TRACE & Ya & Tidak & Ya & Tidak \\
    BPJPH SIHALAL & Ya & Tidak & Ya & Tidak \\
    \textbf{HalalChain} & \textbf{Ya} & \textbf{Ya} & \textbf{Ya} & \textbf{Ya} \\
    \bottomrule
  \end{tabularx}
\end{table}


\subsection{Ketertelusuran Rantai Pasok Blockchain}
Tinjauan sistematis mengidentifikasi ketidakberubahan, keterauditabilitas, dan kepercayaan multi-pemangku kepentingan sebagai motivasi utama blockchain dalam rantai pasok~\cite{casino2021systematic,kshetri2018blockchain,galvez2018future}.
Ledger berizin (Hyperledger Fabric) mendominasi percontohan enterprise karena menawarkan manajemen identitas, saluran privat, dan throughput yang dapat diprediksi~\cite{andoni2019blockchain,abeyratne2016blockchain}.
Rantai publik memungkinkan akses baca tanpa izin---konsumen mana pun dengan pengenal batch dapat memverifikasi status tanpa kerja sama vendor---dengan biaya transaksi dan onboarding dompet bagi penulis~\cite{underwood2016blockchain}.
Rollup Layer~2 seperti Optimistic rollup (termasuk OP Stack yang mendasari Base) mengelompokkan transaksi ke L1 sambil mewarisi asumsi keamanan Ethereum, sehingga menurunkan biaya per operasi~\cite{kalodner2018arbitrum,optimism2024spec}.
HalalChain mengadopsi L2 publik tepat untuk menjaga transparansi baca bagi konsumen sambil menjaga biaya tulis dalam jangkauan UMKM.

Tinjauan industri pangan mencatat pola berulang: hash atau pengenal on-chain, muatan off-chain, dan log peristiwa untuk provenance~\cite{rejeb2020blockchain,kouhizadeh2018blockchain}.
Caro \emph{dkk.} mengimplementasikan ketertelusuran agri-food dengan rantai peristiwa berbasis blockchain~\cite{caro2018blockchain}; HalalChain berbeda dengan menempatkan \emph{status sertifikasi} dan \emph{keputusan auditor} sebagai pusat, bukan hanya peristiwa logistik.

\subsection{Sistem Integritas Halal}
Literatur rantai pasok halal menekankan ketertelusuran dari farm to fork, keaslian sertifikasi, dan pengakuan lintas batas~\cite{sidarto2021halal,rejeb2018halal,rashid2018halal}.
Sidarto dan Hamka mengusulkan blockchain untuk ketertelusuran unggas di Indonesia, menyoroti tantangan integrasi dengan proses BPJPH yang ada~\cite{sidarto2021halal}.
Rejeb \emph{dkk.} menggabungkan HACCP, blockchain, dan IoT untuk rantai pasok daging halal~\cite{rejeb2018halal}.
Usulan terbaru menggabungkan blockchain dengan verifikasi QR dan pemeriksaan dokumen berbantuan AI~\cite{alourani2025halal,marianingsih2024halal}.

Namun, sedikit sistem yang memodelkan \emph{penolakan auditor} sebagai status terminal on-chain dengan jalur pengajuan ulang terstruktur.
Dalam praktik, aplikasi yang ditolak dikoreksi di luar ledger dan masuk kembali ke alur kerja secara tidak transparan---melemahkan akuntabilitas (\emph{Hisba}).
HalalChain mengodekan alasan penolakan, CID dokumen audit, dan tautan \texttt{parentBatchId} sehingga konsumen dan regulator dapat memeriksa jejak keputusan lengkap.

\subsection{Penyimpanan Off-Chain dan IPFS}
Menyimpan sertifikat PDF dan laporan laboratorium lengkap on-chain tidak ekonomis: gas berskala dengan ukuran calldata, dan node rantai mereplikasi semua data secara global~\cite{benet2014ipfs,christidis2016blockchain}.
Penyimpanan berbasis alamat konten melalui IPFS menurunkan pengenal (CID) dari hash dokumen; manipulasi apa pun mengubah CID, sehingga memungkinkan pemeriksaan integritas ketika CID dikaitkan on-chain~\cite{benet2014ipfs}.
Christidis dan Devetsikiotis meninjau pola integrasi kontrak pintar dan IoT di mana blockchain menyimpan referensi, bukan muatan~\cite{christidis2016blockchain}.
HalalChain mengikuti pola ini: kontrak menyimpan string \texttt{ipfsCid} dan \texttt{auditIpfsCid}; Pinata menyediakan pinning agar gateway tetap dapat diakses selama evaluasi.

\subsection{Kesenjangan Penelitian}
Sistem digital halal yang ada sebagian besar mengandalkan basis data terpusat atau blockchain enterprise tanpa verifikasi konsumen tanpa dompet.
HalalChain menggabungkan: (i)~verifikabilitas L2 publik; (ii)~peran produsen/auditor terpisah RBAC; (iii)~riwayat revisi permanen setelah penolakan; dan (iv)~pemetaan eksplisit persyaratan desain etis ke properti protokol (Bagian~\ref{sec:prelim}).
Tolok ukur gas L2 yang dipublikasikan dari penelitian sebelumnya~\cite{kalodner2018arbitrum} menginformasikan motivasi biaya; metrik mainnet atau testnet khusus HalalChain dilaporkan hanya setelah deployment (Bagian~\ref{sec:eval-b}).

\section{Prasyarat: Konteks Kepercayaan Halal dan Persyaratan Desain Etis}
\label{sec:prelim}

Bagian ini memberikan latar kontekstual tentang kesenjangan kepercayaan sertifikasi halal dan menurunkan \emph{persyaratan desain rekayasa} dari prinsip etis Islam.
Bagian ini bukan putusan agama (fatwa); melainkan menginformasikan desain sistem secara analog dengan value-sensitive design, di mana pertimbangan moral membentuk pilihan teknis tanpa menggantikan keahlian domain~\cite{friedman2017value}.

\subsection{Konteks Sertifikasi Halal di Indonesia}
Ekosistem sertifikasi halal Indonesia melibatkan registrasi BPJPH, klasifikasi produk, dan audit MUI untuk banyak kategori~\cite{sidarto2021halal,bpjph2022regulation}.
Produsen mengajukan deklarasi bahan, deskripsi proses, dan dukungan laboratorium atau atestasi pemasok.
Setelah disetujui, label halal dapat diterapkan pada produk atau lini produksi tertentu; pasar ekspor semakin meminta bukti ketertelusuran di luar label itu sendiri~\cite{marianingsih2024halal}.

Masalah yang dilaporkan meliputi pemalsuan sertifikat, pelabelan ulang batch tidak tersertifikasi, dan ketidakmampuan konsumen memverifikasi klaim tingkat batch secara mandiri pada titik pembelian~\cite{rashid2018halal,kshetri2018blockchain}.
Produsen UMKM---produsen pangan skala kecil, industri rumahan yang berkembang, dan koperasi regional---sering tidak memiliki staf TI untuk mengoperasikan suite ketertelusuran enterprise.
HalalChain oleh karena itu menargetkan atestasi \emph{tingkat batch} dengan field on-chain minimal dan halaman verifikasi konsumen ramah seluler.

\subsection{Aktor dan Batas Kepercayaan}
Empat kelas aktor berinteraksi dengan protokol:
\begin{itemize}
  \item \textbf{Produsen (UMKM):} Mendaftarkan batch dan mengunggah bukti ke IPFS; tidak dapat memverifikasi status halal sendiri.
  \item \textbf{Auditor:} Meninjau bukti dan mengubah status batch menjadi Verified, Rejected, atau Revoked; tindakan ditandatangani on-chain.
  \item \textbf{Konsumen:} Membaca status batch publik melalui URL tertaut QR; tidak memegang kunci atau token.
  \item \textbf{Administrator:} Memberikan \texttt{PRODUCER\_ROLE} dan \texttt{AUDITOR\_ROLE}; titik sempit kepercayaan dibahas di Bagian~\ref{sec:discussion}.
\end{itemize}
Protokol \emph{tidak} secara otomatis memvalidasi kimia bahan atau metode penyembelihan; protokol mencatat \emph{auditor mana} menyetujui \emph{dokumen mana} pada \emph{waktu mana}, mengalihkan kepercayaan dari basis data tidak transparan ke catatan yang dapat diperiksa.

\subsection{Dari Prinsip ke Persyaratan Desain}
Prinsip etis Islam menginformasikan \emph{apa} yang harus dioptimalkan sistem, diekspresikan sebagai properti yang dapat diuji.
Kami memilih empat prinsip yang sering dikutip dalam literatur rantai pasok halal dan memetakannya ke mekanisme HalalChain pada Tabel~\ref{tab:ethics}.

\begin{table}[t]
  \caption{Pemetaan persyaratan desain etis ke implementasi HalalChain.}
  \label{tab:ethics}
  \centering
  \footnotesize
  \setlength{\tabcolsep}{3pt}
  \begin{tabularx}{\columnwidth}{@{}>{\raggedright\arraybackslash}p{0.24\columnwidth}>{\raggedright\arraybackslash}p{0.26\columnwidth}X@{}}
    \toprule
    \textbf{Prinsip} & \textbf{Persyaratan} & \textbf{Implementasi} \\
    \midrule
    Amanah (kepercayaan) & Catatan audit permanen & Transisi status tambah-saja; batch ditolak tidak dapat diverifikasi ulang \\
    Thoyyib (kesucian) & Tampilan halal hanya jika terverifikasi & UI konsumen menampilkan sertifikasi hanya jika \texttt{status = Verified} \\
    Adl (keadilan) & Akses data terdistribusi & Baca rantai publik + gateway IPFS; tidak ada gerbang vendor tunggal untuk status \\
    Hisba (akuntabilitas) & Atribusi auditor & Alamat \texttt{auditor}, \texttt{auditedAt}, \texttt{auditIpfsCid} opsional \\
    \bottomrule
  \end{tabularx}
\end{table}

\emph{Amanah} dioperasionalkan oleh penjaga kontrak pintar: setelah \texttt{rejectBatch} dieksekusi, \texttt{verifyBatch} pada \texttt{batchId} yang sama gagal dengan \texttt{InvalidStatus}, mencegah ``persetujuan'' retroaktif atas batch yang gagal.
\emph{Thoyyib} muncul di dasbor konsumen, yang menampilkan indikator halal positif hanya untuk batch Verified---status Pending dan Rejected menampilkan pesan tidak tersertifikasi eksplisit dalam bahasa Inggris dan Indonesia.
\emph{Adl} sebagian terpenuhi oleh akses baca publik tetapi dibatasi karena penugasan peran tetap terpusat di bawah \texttt{DEFAULT\_ADMIN\_ROLE}.
\emph{Hisba} mengikat setiap keputusan ke alamat Ethereum dan cap waktu blok, memungkinkan akuntabilitas pasca-hoc jika kunci auditor diketahui.

Prototipe saat ini memberikan \texttt{AUDITOR\_ROLE} melalui administrator (OpenZeppelin \texttt{DEFAULT\_ADMIN\_ROLE}), yang sebagian membatasi \emph{Adl}; Bagian~\ref{sec:discussion} membahas mitigasi multi-tanda tangan dan berbasis DAO~\cite{gnosis2024safe}.

\FloatBarrier

\section{Desain Sistem}
\label{sec:design}

\subsection{Gambaran Arsitektur}
Gambar~\ref{fig:arch} mengilustrasikan arsitektur HalalChain melintasi penyimpanan dokumen off-chain dan pengaitan status on-chain.
Produsen mengunggah dokumen pendukung ke IPFS melalui integrasi Pinata sisi server, memperoleh CID, dan memanggil \texttt{registerBatch} pada Base L2 dengan CID tersebut.
Auditor mengambil dokumen melalui gateway IPFS, mencatat keputusan on-chain, dan secara opsional melampirkan CID laporan audit.
Konsumen mengkueri status batch secara read-only melalui \texttt{getBatch} via endpoint JSON-RPC yang disematkan dalam URL web; kode QR mengodekan URL ini sehingga pemindaian tidak memerlukan perangkat lunak dompet.

\begin{figure}[!t]
  \centering
  \resizebox{\columnwidth}{!}{%
  \begin{tikzpicture}[
    font=\footnotesize,
    >=Stealth,
    actor/.style={
      draw,
      rounded corners=2pt,
      line width=0.4pt,
      fill=gray!6,
      minimum width=2.05cm,
      minimum height=0.9cm,
      align=center,
      inner sep=4pt
    },
    core/.style={
      draw,
      rounded corners=2pt,
      line width=0.5pt,
      fill=gray!14,
      minimum width=2.35cm,
      minimum height=0.95cm,
      align=center,
      inner sep=4pt
    },
    store/.style={
      draw,
      rounded corners=2pt,
      line width=0.4pt,
      fill=gray!10,
      minimum width=2.05cm,
      minimum height=0.9cm,
      align=center,
      inner sep=4pt
    },
    lbl/.style={
      font=\scriptsize,
      fill=white,
      inner sep=1.5pt,
      text=black
    }
  ]
    \node[actor] (producer) {Produsen\\UMKM};
    \node[store, right=1.25cm of producer] (ipfs) {IPFS\\Pinata};
    \node[core] at ($(producer)!0.5!(ipfs) + (0,-1.45)$) (chain) {Base L2\\HalalChain.sol};

    \node[actor, below left=1.05cm and 0.55cm of chain] (consumer) {Konsumen\\verifikasi QR};
    \node[actor, below right=1.05cm and 0.55cm of chain] (auditor) {Auditor};

    \draw[->, thick] (producer.east) -- node[lbl, above]{unggah dok.} (ipfs.west);
    \draw[->, thick] (producer.south) -- ++(0,-0.55) -|
      node[lbl, pos=0.2, below]{registerBatch} (chain.west);

    \draw[->, thick, dashed] (ipfs.south) -- node[lbl, right]{simpan CID} (chain.north);

    \draw[->, thick] (chain.south west) -- node[lbl, left, pos=0.35]{baca status} (consumer.north);
    \draw[->, thick] (auditor.north) -- node[lbl, right, pos=0.35]{verifikasi / tolak} (chain.south east);

  \end{tikzpicture}%
  }
  \caption{Arsitektur HalalChain: dokumen pada IPFS, jangkar dan status pada Base L2.}
  \label{fig:arch}
\end{figure}


Desain memisahkan \emph{bukti} (hosting besar dan dapat berubah pada IPFS) dari \emph{komitmen} (kecil, direplikasi, permanen pada L2).
Tabel~\ref{tab:datasplit} merangkum pemisahan tersebut.

\begin{table}[t]
  \caption{Penempatan data on-chain vs.\ off-chain.}
  \label{tab:datasplit}
  \centering
  \footnotesize
  \begin{tabularx}{\columnwidth}{@{}>{\raggedright\arraybackslash}p{0.36\columnwidth}>{\raggedright\arraybackslash}p{0.18\columnwidth}X@{}}
    \toprule
    \textbf{Data} & \textbf{Lokasi} & \textbf{Alasan} \\
    \midrule
    Status batch, cap waktu & On-chain & FSM tahan manipulasi \\
    Alamat produsen/auditor & On-chain & Akuntabilitas \\
    CID IPFS & On-chain & Jangkar berbasis alamat konten \\
    Alasan penolakan (string) & On-chain & Jejak audit permanen \\
    PDF, gambar, laporan lab & IPFS & Biaya dan ukuran \\
    Berkas pendukung audit & IPFS & Sama seperti dokumen produsen \\
    \bottomrule
  \end{tabularx}
\end{table}

\subsection{Siklus Hidup Batch}
Status batch mengikuti mesin status terbatas (FSM) dengan lima nilai enumerasi: \texttt{Unknown}, \texttt{Pending}, \texttt{Verified}, \texttt{Rejected}, dan \texttt{Revoked}.
Tabel~\ref{tab:fsm} mencantumkan transisi yang diizinkan; semua transisi lain gagal pada lapisan kontrak.

\begin{table}[t]
  \caption{Transisi status batch yang ditegakkan oleh \texttt{HalalChain.sol}.}
  \label{tab:fsm}
  \centering
  \footnotesize
  \setlength{\tabcolsep}{3pt}
  \begin{tabularx}{\columnwidth}{@{}>{\raggedright\arraybackslash}p{0.14\columnwidth}>{\raggedright\arraybackslash}p{0.20\columnwidth}>{\raggedright\arraybackslash}X@{}}
    \toprule
    \textbf{Dari} & \textbf{Ke} & \textbf{Pemicu} \\
    \midrule
    --- & Pending & \makecell[tl]{\texttt{registerBatch},\\\texttt{registerRevision}} \\
    Pending & Verified & \makecell[tl]{\texttt{verifyBatch}\\\textit{(auditor)}} \\
    Pending & Rejected & \makecell[tl]{\texttt{rejectBatch}\\\textit{(auditor)}} \\
    Verified & Revoked & \makecell[tl]{\texttt{revokeBatch}\\\textit{(auditor)}} \\
    Rejected & Pending (id baru) & \texttt{registerRevision} \\
    Rejected & Verified & \emph{Dilarang} (gagal) \\
    \bottomrule
  \end{tabularx}
\end{table}

Batch yang ditolak tidak dapat diverifikasi ulang pada tempatnya; produsen mengajukan \emph{revisi} sebagai batch baru yang ditautkan melalui \texttt{parentBatchId}, mempertahankan alasan penolakan asli dan metadata audit untuk pemeriksaan.
Pencabutan hanya berlaku untuk batch yang sebelumnya Verified---misalnya, ketika pengawasan pasca-pasar menemukan kecurangan---dan mencatat string alasan on-chain.

\subsection{Spesifikasi Kontrak Pintar}
Kontrak \texttt{HalalChain.sol} (Solidity ${\hat{}}0.8.20$, OpenZeppelin \texttt{AccessControl}~\cite{openzeppelin2024access}) mendefinisikan konstanta peran \texttt{PRODUCER\_ROLE}, \texttt{AUDITOR\_ROLE}, dan \texttt{DEFAULT\_ADMIN\_ROLE}.
Kesalahan kustom (\texttt{InvalidBatch}, \texttt{InvalidStatus}, \texttt{NotOriginalProducer}, \texttt{ParentNotRejected}) menggantikan string revert untuk menghemat gas dan meningkatkan kejelasan integrator.

\textbf{Field struct Batch:}
\texttt{producer}, \texttt{productName}, \texttt{ipfsCid}, \texttt{createdAt}, \texttt{status}, \texttt{auditor}, \texttt{auditedAt}, \texttt{auditIpfsCid}, \texttt{rejectReason}, \texttt{parentBatchId}.

\textbf{Fungsi utama:}
\begin{itemize}
  \item \texttt{registerBatch(productName, ipfsCid)} --- membuat batch dengan \texttt{parentBatchId = 0}, status Pending; memancarkan \texttt{BatchRegistered}.
  \item \texttt{registerRevision(parentBatchId, ipfsCid)} --- mensyaratkan induk Rejected dan \texttt{msg.sender == parent.producer}; mewarisi \texttt{productName}; memancarkan peristiwa dengan tautan induk.
  \item \texttt{verifyBatch(batchId, auditIpfsCid)} --- Pending $\rightarrow$ Verified; mengatur field auditor; memancarkan \texttt{BatchVerified}.
  \item \texttt{rejectBatch(batchId, reason, auditIpfsCid)} --- Pending $\rightarrow$ Rejected; menyimpan alasan; memancarkan \texttt{BatchRejected}.
  \item \texttt{revokeBatch(batchId, reason)} --- Verified $\rightarrow$ Revoked; memancarkan \texttt{BatchRevoked}.
  \item \texttt{getBatch(batchId)} --- view publik mengembalikan struct lengkap; digunakan halaman konsumen tanpa biaya transaksi.
\end{itemize}

Peristiwa diindeks pada \texttt{batchId} dan alamat peserta jika berlaku, memungkinkan block explorer dan pengindeks subgraph merekonstruksi linimasa audit tanpa API proprietari.

\subsection{Penyimpanan Off-Chain dan Integritas}
Dokumen dipin pada IPFS melalui Pinata; hanya CID yang disimpan on-chain~\cite{benet2014ipfs,christidis2016blockchain}.
Integritas mengikuti penandaan alamat konten: jika gateway mengembalikan byte yang hash-nya tidak cocok dengan CID yang dikaitkan, klien harus memperlakukan dokumen sebagai tidak terpercaya.
Ketersediaan bergantung pada kebijakan pinning---CID yang tidak dipin dapat menjadi tidak dapat dijangkau meskipun referensi on-chain tetap ada---keterbatasan yang dianalisis di Bagian~\ref{sec:discussion}.
API unggah tidak pernah mengekspos kredensial JWT Pinata ke browser; hanya rute Next.js sisi server yang menyimpan rahasia.

\subsection{Jalur Verifikasi Konsumen}
Konsumen mengakses \texttt{/verify/\{batchId\}} pada frontend yang di-deploy.
Halaman memanggil \texttt{getBatch} melalui endpoint RPC publik (klien baca Viem/Wagmi), merender UI spesifik status (Verified / Pending / Rejected / Revoked), menampilkan tautan IPFS untuk dokumen produsen dan audit, serta menampilkan \texttt{parentBatchId} jika ada.
Tidak ada materi kunci privat yang dimuat pada rute konsumen, memenuhi persyaratan tanpa dompet pada Tabel~\ref{tab:related}.

\section{Implementasi}
\label{sec:implementation}

Tabel~\ref{tab:stack} merangkum tumpukan teknologi.
Repositori sumber terbuka tersedia di \url{https://github.com/Fatihmaull/halal-chain}.

\begin{table}[t]
  \caption{Tumpukan teknologi implementasi HalalChain.}
  \label{tab:stack}
  \centering
  \footnotesize
  \setlength{\tabcolsep}{3pt}
  \begin{tabularx}{\columnwidth}{@{}>{\raggedright\arraybackslash}p{0.22\columnwidth}X@{}}
    \toprule
    \textbf{Lapisan} & \textbf{Teknologi} \\
    \midrule
    Blockchain & Solidity 0.8.20, OpenZeppelin 5.x, Base Sepolia (84532) \\
    Pengembangan & Hardhat 3.7, Mocha/Chai (7 pengujian), skrip TypeScript \\
    Frontend & Next.js 16, Wagmi v3, Viem 2.x, Tailwind CSS \\
    Penyimpanan & IPFS via Pinata (\texttt{POST /api/ipfs/upload}) \\
    Deployment & \makecell[tl]{Hardhat (\texttt{deploy:local},\\\texttt{deploy:base-sepolia}), Vercel} \\
    \bottomrule
  \end{tabularx}
\end{table}

\subsection{Kontrak Pintar}
Kontrak \texttt{HalalChain} memperluas OpenZeppelin \texttt{AccessControl}~\cite{openzeppelin2024access}, menggantikan pemetaan auditor pemilik tunggal sebelumnya dengan peran eksplisit yang dapat diberikan oleh \texttt{DEFAULT\_ADMIN\_ROLE}.
Konstruktor menetapkan admin ke deployer; skrip deployment secara opsional memberikan peran produsen dan auditor ke alamat yang dikonfigurasi untuk demo testnet.

Pengujian unit (\texttt{contracts/test/HalalChain.test.ts}) mencakup:
\begin{itemize}
  \item pendaftaran batch dan penugasan \texttt{nextBatchId} monoton;
  \item verifikasi dan penolakan dengan penegakan peran;
  \item revisi setelah penolakan dengan penjaga \texttt{NotOriginalProducer} dan \texttt{ParentNotRejected};
  \item pencabutan dari status Verified;
  \item panggilan tidak berwenang dari alamat tanpa \texttt{PRODUCER\_ROLE} atau \texttt{AUDITOR\_ROLE};
  \item pencarian \texttt{batchId} tidak valid yang gagal dengan \texttt{InvalidBatch}.
\end{itemize}

Skrip Hardhat meliputi \texttt{deploy.ts} (lokal atau Base Sepolia), \texttt{export-abi.ts} (salinan ABI ke \texttt{web/src/contracts/}), dan \texttt{evaluate-gas.ts} (agregasi gas receipt untuk skenario evaluasi).

\subsection{Aplikasi Web}
Tiga dasbor spesifik peran berbagi tata letak umum dan string dwibahasa (Inggris/Indonesia):

\textbf{Dasbor produsen} memandu pengguna melalui penamaan produk, unggah multi-berkas, pinning IPFS, dan pengiriman \texttt{registerBatch}.
Setelah konfirmasi, UI menampilkan \texttt{batchId} yang ditetapkan, hash transaksi, dan kode QR yang mengarah ke URL verifikasi konsumen.
Riwayat batch mencantumkan registrasi sebelumnya dengan lencana status dan tautan ke batch induk jika berlaku.

\textbf{Dasbor auditor} menyajikan antrean tertunda yang diambil melalui iterasi \texttt{getBatch} atau log peristiwa.
Auditor membuka tautan dokumen IPFS di tab baru, lalu mengirim \texttt{verifyBatch}, \texttt{rejectBatch} (dengan teks alasan), atau \texttt{revokeBatch} melalui tulisan terhubung dompet.
Modal konfirmasi mengurangi perubahan status yang tidak disengaja.

\textbf{Halaman verifikasi konsumen} (\texttt{/verify/[batchId]}) melakukan panggilan RPC read-only.
Halaman merender hijau ``Halal Certified'' / ``Tersertifikasi Halal'' untuk batch Verified, pesan penolakan merah dengan \texttt{rejectReason} untuk batch Rejected, dan status pending/revoked netral untuk sisanya.
Tautan batch induk memungkinkan ketertelusuran melintasi rantai revisi.

\subsection{Alur Unggah IPFS}
Rute API Next.js (\texttt{POST /api/ipfs/upload}) menerima unggah berkas \texttt{multipart/form-data}, mem-pin melalui JWT Pinata sisi server, dan mengembalikan muatan JSON CID ke klien.
Variabel lingkungan \texttt{PINATA\_JWT} diperlukan untuk pinning produksi; ketika tidak ada, mode mock pengembangan menghasilkan CID lokal deterministik agar alur UI tetap dapat diuji tanpa layanan eksternal.
Tolok ukur yang dimaksudkan untuk publikasi (Bagian~\ref{sec:eval-b}) hanya akan menggunakan pin Pinata nyata.

\subsection{Konfigurasi dan Deployment}
Variabel lingkungan frontend meliputi \texttt{NEXT\_PUBLIC\_HALALCHAIN\_ADDRESS}, \texttt{NEXT\_PUBLIC\_CHAIN\_ID}, \texttt{NEXT\_PUBLIC\_SEPOLIA\_RPC\_URL}, dan \texttt{NEXT\_PUBLIC\_DEMO\_URL} untuk URL dasar QR.
Dasbor konsumen di-deploy pada Vercel (\url{https://web-six-ivory-36.vercel.app/}); deployment kontrak pada Ethereum Sepolia atau Base Sepolia memerlukan \texttt{DEPLOYER\_PRIVATE\_KEY} dan ETH dari faucet (Lampiran~\ref{app:repro}).

\section{Evaluasi}
\label{sec:eval}

Kami menyusun evaluasi dalam dua tingkat selaras dengan kematangan deployment: (VI-A) analisis kualitatif dan validasi fungsional yang tersedia pada jaringan Hardhat lokal saat ini; dan (VI-B) tolok ukur kuantitatif Base Sepolia yang menunggu deployment testnet berdana.
Pemisahan ini menghindari pencampuran pengukuran mesin pengembang dengan biaya jaringan publik~\cite{casino2021systematic}.

\subsection{Analisis Kualitatif dan Validasi Fungsional}
\label{sec:eval-a}

\subsubsection{Properti Keamanan}
Tabel~\ref{tab:security} memetakan persyaratan desain dari Bagian~\ref{sec:prelim} ke perilaku protokol yang dapat diamati, divalidasi oleh pengujian unit dan eksekusi skenario manual.

\begin{table}[t]
  \caption{Properti keamanan dan metode validasi.}
  \label{tab:security}
  \centering
  \footnotesize
  \setlength{\tabcolsep}{3pt}
  \begin{tabularx}{\columnwidth}{@{}>{\raggedright\arraybackslash}p{0.34\columnwidth}X@{}}
    \toprule
    \textbf{Properti} & \textbf{Validasi} \\
    \midrule
    Ketidakberubahan penolakan & \texttt{verifyBatch} pada Rejected gagal \\
    Ketertelusuran revisi & \texttt{parentBatchId} dipertahankan; alasan induk dapat dibaca \\
    Pemisahan peran & Non-produsen tidak dapat mendaftar; non-auditor tidak dapat memverifikasi \\
    Identitas batch & \texttt{batchId} tidak valid gagal pada semua mutator \\
    Integritas pencabutan & Hanya Verified $\rightarrow$ Revoked diizinkan \\
    \bottomrule
  \end{tabularx}
\end{table}

\paragraph{Ketidakberubahan.}
Setelah \texttt{rejectBatch} dieksekusi, \texttt{verifyBatch} pada pengenal yang sama gagal dengan \texttt{InvalidStatus}.
Skenario~C (di bawah) mengonfirmasi bahwa upaya verifikasi berulang tidak mengubah status on-chain---mendukung \emph{Amanah}.

\paragraph{Ketertelusuran revisi.}
\texttt{registerRevision} membuat \texttt{batchId} baru sambil menyimpan \texttt{parentBatchId}.
Konsumen yang memeriksa revisi terverifikasi dapat menavigasi ke induk yang ditolak dan membaca \texttt{rejectReason} serta \texttt{ipfsCid} aslinya, mendukung transparansi audit.

\paragraph{Kontrol akses.}
Modifier OpenZeppelin \texttt{onlyRole} mengawal fungsi mutasi.
Pengujian memastikan alamat tidak berwenang tidak dapat mendaftar, memverifikasi, menolak, atau mencabut, mengatasi kekhawatiran eskalasi hak istimewa dasar pada prototipe v1.

\subsubsection{Skenario Fungsional}
Tiga skenario end-to-end dari \texttt{docs/EVALUATION.md} dieksekusi pada testnet Ethereum Sepolia (Lampiran~\ref{app:repro}) pada 2026-06-08, dengan bukti diarsipkan di \path{docs/evaluation/}.

\paragraph{Skenario A (jalur sukses).}
Skenario~A (batch \#1): pendaftaran produsen dan verifikasi auditor pada testnet Ethereum Sepolia; halaman konsumen di \url{https://web-six-ivory-36.vercel.app/verify/1} menampilkan status terverifikasi.

\paragraph{Skenario B (penolakan dan revisi).}
Skenario~B: batch \#2 ditolak (\url{https://web-six-ivory-36.vercel.app/verify/2}), revisi batch \#3 terverifikasi (\url{https://web-six-ivory-36.vercel.app/verify/3}) dengan tautan induk dipertahankan.

\paragraph{Skenario C (ketahanan manipulasi).}
Skenario~C: \texttt{verifyBatch} pada batch ditolak \#2 gagal (\texttt{InvalidStatus}); status on-chain tidak berubah.

Validasi fungsional pada Ethereum Sepolia mengonfirmasi kebenaran protokol dan baca konsumen tanpa dompet melalui frontend Next.js; biaya gas pada testnet L1 membatasi atas kelayakan deployment L2~\cite{optimism2024spec}.


\subsubsection{Observasi Gas Jaringan Pengembangan}
Receipt Hardhat localhost menyediakan penggunaan gas \emph{awal} untuk penentuan ukuran kontrak (Tabel~\ref{tab:gas-local}).
Angka ini \textbf{bukan} dilaporkan sebagai biaya Base Sepolia---harga opcode dan biaya data L1 berbeda pada rollup~\cite{optimism2024spec}.

\begin{table}[t]
  \caption{Gas awal pada Hardhat localhost (hanya pengembangan; bukan Base Sepolia).}
  \label{tab:gas-local}
  \centering
  \footnotesize
  \begin{tabularx}{\columnwidth}{@{}Xr@{}}
    \toprule
    \textbf{Operasi} & \textbf{Gas Digunakan} \\
    \midrule
    \texttt{registerBatch} & 155{,}506 \\
    \texttt{verifyBatch} & 79{,}214 \\
    \texttt{rejectBatch} & 103{,}533 \\
    \texttt{registerRevision} & 181{,}835 \\
    \texttt{getBatch} (view) & 0 \\
    \bottomrule
  \end{tabularx}
\end{table}

\subsection{Tolok Ukur Kuantitatif Testnet}
\label{sec:eval-b}

Tabel~\ref{tab:gas} melaporkan penggunaan gas dan biaya transaksi yang diukur pada Ethereum Sepolia (chain ID 11155111) pada 2026-06-08.
Kontrak tolok ukur gas: \url{https://sepolia.etherscan.io/address/0x0Ccf5d905787E5E7EA698469fAe38a1517Ff8c29}.
Semua nilai diambil dari receipt on-chain melalui \texttt{npm run evaluate:gas:eth-sepolia} dan diverifikasi silang pada Etherscan.

\begin{table}[t]
  \caption{Metrik gas dan biaya pada Ethereum Sepolia (2026-06-08).}
  \label{tab:gas}
  \centering
  \footnotesize
  \begin{tabularx}{\columnwidth}{@{}Xrr@{}}
    \toprule
    \textbf{Operasi} & \textbf{Gas Digunakan} & \textbf{Biaya (ETH)} \\
    \midrule
    \texttt{registerBatch} & 155,506 & 0,00078367 \\
    \texttt{verifyBatch} & 79,214 & 0,00034832 \\
    \texttt{rejectBatch} & 103,533 & 0,00055169 \\
    \texttt{registerRevision} & 181,835 & 0,00096602 \\
    \texttt{getBatch} (baca RPC) & 0 & 0 \\
    \bottomrule
  \end{tabularx}
\end{table}

Frontend konsumen di-deploy pada Vercel di \url{https://web-six-ivory-36.vercel.app/}; metrik yang direncanakan meliputi waktu konfirmasi blok (pendaftaran $\rightarrow$ verifikasi), durasi unggah IPFS via Pinata, dan time-to-first-byte halaman konsumen pada host ini.
Jalankan \texttt{npm run evaluate:gas} di \texttt{contracts/} setelah deployment Base Sepolia (Lampiran~\ref{app:repro}).

\section{Diskusi}
\label{sec:discussion}

\subsection{Asumsi Kepercayaan dan Keterbatasan}
Kami menyatakan secara eksplisit asumsi kepercayaan berikut agar pengadopsi dapat menalar risiko residual:

\begin{enumerate}
  \item \textbf{Kejujuran administrator} --- \texttt{DEFAULT\_ADMIN\_ROLE} memberikan hak istimewa auditor dan produsen.
  Kompromi kunci admin atau kolusi dengan auditor jahat merusak integritas sertifikasi meskipun log permanen.
  \item \textbf{Ketelitian auditor} --- catatan on-chain membuktikan bahwa alamat auditor menyetujui batch pada waktu tertentu; catatan tidak secara kriptografis membuktikan kepatuhan halal bahan, metode penyembelihan, atau perilaku pemasok.
  HalalChain adalah \emph{jejak audit}, bukan klasifikator halal otomatis.
  \item \textbf{Ketersediaan pinning IPFS} --- dokumen memerlukan pin Pinata (atau self-hosted) yang aktif.
  CID yang tidak dipin atau kedaluwarsa menjadi tidak tersedia off-chain sementara CID on-chain tetap ada, menghasilkan status referensi ``menggantung'' yang harus ditafsirkan konsumen dengan hati-hati.
  \item \textbf{Kelangsungan sequencer L2} --- Base mengandalkan infrastruktur rollup; gangguan berkepanjangan menunda tulisan tetapi tidak mengubah status yang dikomitkan secara diam-diam.
  \item \textbf{Ruang lingkup testnet} --- evaluasi pada Base Sepolia tidak menyiratkan integrasi BPJPH produksi, kesetaraan hukum dengan registrasi SIHALAL, atau kesiapan mainnet tanpa tinjauan tata kelola dan manajemen kunci lebih lanjut.
\end{enumerate}

Keterbatasan ini sebagian membatasi persyaratan desain \emph{Adl} (keadilan) hingga tata kelola auditor multi-tanda tangan atau penugasan peran berbasis DAO diimplementasikan~\cite{gnosis2024safe}.
Versi mendatang dapat mensyaratkan persetujuan auditor \texttt{m}-dari-\texttt{n} untuk verifikasi, menyelaraskan kebijakan on-chain dengan lembaga halal institusional.

\subsection{Perbandingan dengan Alternatif Terpusat}
Basis data terpusat menawarkan kompleksitas operasional lebih rendah, alat SQL yang familiar, dan integrasi alur kerja BPJPH yang mudah~\cite{underwood2016blockchain,viriyasitavat2019blockchain}.
Administrator dapat mencabut akses, memperbaiki kesalahan ketik, dan menjalankan laporan batch tanpa biaya gas.
Namun, model terpusat memungkinkan alterasi catatan retroaktif oleh pengguna berhak dan tidak menyediakan verifikasi pihak ketiga tanpa izin---konsumen harus mempercayai respons API operator.

HalalChain menukar properti ini dengan verifikabilitas publik dan jejak audit tambah-saja.
Produsen dan auditor harus mengelola dompet Ethereum dan ETH (atau ETH yang di-bridge pada Base) untuk tulisan; pinning IPFS menimbulkan biaya berulang atau per-GB.
Bagi produsen UMKM, onboarding bersubsidi (meta-transaksi, paymaster, atau sponsor institusional) mungkin diperlukan untuk adopsi skala besar.

Dibandingkan blockchain enterprise (Hyperledger), HalalChain lebih mengutamakan transparansi daripada privasi transaksi: status batch dapat dibaca dunia.
Sertifikasi halal pada dasarnya adalah \emph{klaim publik}; saluran privat mungkin penting untuk formulasi rahasia dagang tetapi di luar ruang lingkup v1.

\subsection{Implikasi bagi UMKM dan Kebijakan}
Mandat pelabelan halal Indonesia menciptakan permintaan alat ketertelusuran terjangkau.
HalalChain menunjukkan arsitektur referensi yang dapat dipilot oleh lembaga terakreditasi BPJPH, koperasi regional, atau tim KKN universitas tanpa melisensikan platform proprietari.
Integrasi dengan SIHALAL memerlukan akses API resmi, pemetaan hukum status on-chain ke nomor registrasi, dan alur kerja akreditasi auditor---tidak ada yang diimplementasikan pada prototipe.

Pasar ekspor semakin mengharapkan ketertelusuran digital~\cite{marianingsih2024halal}.
Halaman verifikasi publik tertaut QR dapat melengkapi label fisik, meskipun penerimaan regulasi bervariasi menurut negara tujuan.

\subsection{Penelitian Lanjutan}
Arah masa depan meliputi:
\begin{itemize}
  \item verifikasi auditor multi-tanda tangan dan pencabutan terkunci waktu;
  \item pemilihan auditor berbasis DAO dengan modul stake atau reputasi;
  \item deployment Base mainnet dengan dasbor pemantauan biaya;
  \item pengaitan sensor IoT (suhu, GPS) melalui jaringan oracle;
  \item integrasi dengan BPJPH SIHALAL jika API tersedia;
  \item verifikasi formal transisi FSM di \texttt{HalalChain.sol};
  \item studi pengguna dengan produsen UMKM dan konsumen tentang kegunaan onboarding dompet dan verifikasi QR.
\end{itemize}

\section{Kesimpulan}
\label{sec:conclusion}

Makalah ini memperkenalkan HalalChain, protokol ketertelusuran batch untuk sertifikasi halal yang menggabungkan catatan audit permanen on-chain pada Base L2 dengan penyimpanan dokumen IPFS off-chain.
Kami mengatasi defisit kepercayaan akibat pemalsuan sertifikat dan manipulasi terpusat dengan mengodekan pendaftaran produsen, keputusan auditor, dan riwayat revisi sebagai transisi status terbatas eksplisit yang ditegakkan oleh kontrak pintar Solidity dengan RBAC OpenZeppelin.

Kontribusi meliputi: (i)~arsitektur tiga peran minimal yang memungkinkan verifikasi konsumen tanpa dompet melalui URL tertaut QR; (ii)~penolakan permanen dan revisi terstruktur melalui tautan \texttt{parentBatchId}; (iii)~pemetaan persyaratan desain etis selaras Maqasid ke properti protokol yang dapat diuji; dan (iv)~implementasi sumber terbuka dengan pengujian unit, skenario fungsional, dan pipeline evaluasi yang dapat direproduksi.

Evaluasi kualitatif mengonfirmasi ketahanan manipulasi pada batch ditolak, kontrol akses berbasis peran, dan integrasi dasbor end-to-end pada jaringan lokal.
Pengukuran gas localhost awal menunjukkan gas di bawah 200k untuk operasi tulis utama, mengisyaratkan kelayakan deployment L2 bagi produsen UMKM, meskipun tolok ukur Base Sepolia masih menunggu deployment berdana.

HalalChain bukan pengganti penetapan halal BPJPH/MUI; HalalChain menyediakan bukti transparan tentang \emph{auditor mana} menyetujui \emph{dokumen batch mana} pada \emph{waktu mana}.
Penelitian lanjutan akan mengatasi tata kelola auditor terdesentralisasi, optimasi biaya mainnet, dan pilot berorientasi kebijakan dengan pemangku kepentingan sertifikasi halal Indonesia.


\section*{Kontribusi Penulis}
Muhammad Fatih Maulana, Rio Ferdana Sudrajat, Muhammad Rahardian Baihaqi, Muhammad Fahriza Pratama, Yazid Zinky Arisona, dan Muhammad Fauzan Lubada berkontribusi pada konseptualisasi, pengembangan perangkat lunak, validasi, dan penulisan.
Seluruh penulis telah membaca dan menyetujui versi terbitan naskah ini.

\section*{Pernyataan Ketersediaan Data}
Kode sumber, kontrak pintar, dan skrip evaluasi tersedia di \url{https://github.com/Fatihmaull/halal-chain}.
Demonstrasi langsung untuk konsumen di-deploy di \url{https://web-six-ivory-36.vercel.app/}.

\section*{Konflik Kepentingan}
Penulis menyatakan tidak ada konflik kepentingan.

\appendices
\section{Reproduksibilitas}
\label{app:repro}

\subsection{Versi Perangkat Lunak}
\begin{itemize}
  \item Solidity 0.8.20, Hardhat 3.7, OpenZeppelin Contracts 5.x
  \item Next.js 16.2, Wagmi 3.6, Viem 2.52
  \item Node.js 20+ direkomendasikan; npm 10+
\end{itemize}

\subsection{Tata Letak Repositori}
Repositori diorganisasi sebagai berikut (GitHub: \texttt{halal-chain}):

\vspace{0.2em}
\noindent\texttt{contracts/}\\
\noindent\hspace{1.1em}\hangindent=1.1em\hangafter=0
\texttt{HalalChain.sol}, pengujian unit, skrip deploy dan gas

\vspace{0.15em}
\noindent\texttt{web/}\\
\noindent\hspace{1.1em}\hangindent=1.1em\hangafter=0
Dasbor Next.js, API unggah IPFS, ABI kontrak

\vspace{0.15em}
\noindent\texttt{docs/}\\
\noindent\hspace{1.1em}\hangindent=1.1em\hangafter=0
\texttt{EVALUATION.md}, \texttt{BASELINE.md}, metrik localhost

\vspace{0.15em}
\noindent\texttt{paper/}\\
\noindent\hspace{1.1em}\hangindent=1.1em\hangafter=0
Sumber IEEEtran untuk naskah ini (Inggris: \texttt{main.tex}; Indonesia: \texttt{main-indo.tex})

\subsection{Build dan Pengujian}
\begin{Code}
cd contracts && npm install && npm test
npm run node
npm run deploy:local
cd ../web && npm install && npm run dev
\end{Code}
Jalankan \texttt{node} dan \texttt{deploy:local} di terminal terpisah.
Tujuh pengujian unit harus lulus.
Akun Hardhat default \#0 menerima \texttt{DEFAULT\_ADMIN\_ROLE} saat deploy; skrip memberikan peran produsen dan auditor ke akun \#1 dan \#2 untuk demo lokal.

\subsection{Deployment Ethereum Sepolia}
Konfigurasi \texttt{contracts/.env}:
\begin{Code}
DEPLOYER_PRIVATE_KEY=0x...
PRODUCER_ADDRESS=0x...
AUDITOR_ADDRESS=0x...
BASE_SEPOLIA_RPC_URL=
  https://sepolia.base.org
\end{Code}
Kemudian jalankan \texttt{npm run deploy:base-sepolia} dan \texttt{npm run evaluate:gas}.
Danai deployer melalui faucet Base Sepolia sebelum menyiarkan transaksi.

\subsection{Catatan Deployment}
\begin{itemize}
  \item Komit Git:\\
  \url{https://github.com/Fatihmaull/halal-chain/commit/313f06c4ee329ac58a8c17eb3399802e077361be}
  \item Kontrak (Ethereum Sepolia):\\
  \url{https://sepolia.etherscan.io/address/0xDaCA688e86F438A7cD6B0C9B69606C67CE85Dc92}
  \item Chain ID: 11155111 (Ethereum Sepolia)
  \item Cap waktu evaluasi: 2026-06-08T12:37:38.107Z
  \item Bundel bukti: \path{docs/evaluation/EVALUATION_REPORT.json}
  \item Demo frontend:\\
  \url{https://web-six-ivory-36.vercel.app/}
  \item Verifikasi batch \#1 (terverifikasi):\\
  \url{https://web-six-ivory-36.vercel.app/verify/1}
  \item Verifikasi batch \#2 (ditolak):\\
  \url{https://web-six-ivory-36.vercel.app/verify/2}
  \item Verifikasi batch \#3 (revisi):\\
  \url{https://web-six-ivory-36.vercel.app/verify/3}
  \item Tutorial setup: \path{docs/ETH_SEPOLIA_SETUP.md}
\end{itemize}

Lihat \texttt{docs/BASELINE.md} untuk alamat dompet default Hardhat yang digunakan dalam pengembangan lokal.


\bibliographystyle{IEEEtran}
\bibliography{references}

\end{document}
